\section{Trying to break it}
\label{sec:break}

Before the interval proofs we tried to break the claim and the lines it rests on.
For every search we report the families, planted adversaries, and deliberately wrong versions that the same search does break.
Everything in this section is \numst.

\subsection{The lines against the lemmas at a strong ripple}
The audit of~\cite{I} (Sec.~VIII and App.~B there) compares every lemma piece with its line on families of states: uniform, edges, corners, several edges at once, large $b$, with the brightness at its ceiling for the corner family, and a grid of $a$ with its end points.
We ran it with the lemma at the ripple $\eta\in\{1/\ell,0.05,0.15,0.25\}$, the target $\gamma(1-\eta)/(1+\eta)$, and the lines of this paper at $\lambda=1/(1+\eta)$, in both forms (remainders at $\bar\eta$, and the chord), for $\ell\in\{55.26,960.8,10^4,10^6,10^{12}\}$: $94\,512$ points, no group and no option total violated (tolerance $10^{-30}$ relative).
Wrong versions, each at $\eta=\tfrac14$ and $\ell\in\{55.26,960.8\}$: the remainders at $1/\ell$ (``$\eta$ forgotten''), $5\,554$--$7\,200$ violations; the target not scaled, $1\,886$--$2\,137$; P2z$\times1.01$, $485$--$646$; N1$\times0.99$, $511$--$687$; B2c$\times0.99$, $470$--$474$; the five lines of FrB1$\times0.70$, $416$--$441$.
The old constant $X_1$ in B1m is \emph{not} seen by this audit, because the group FrB1 has slack; it breaks the written step of the chain, which the next audit checks.

The 30 groups of formulas of~\cite{I} (lemma piece, closed form, elementary steps, line) were evaluated at $\eta=\tfrac14$ for $\ell\in\{55.26,955.8,10^6,10^{12}\}$.
With the steps of~\cite{I}, only FrB1 breaks, in its first step; with $X_{1,\rm B1}$ no group breaks.
Wrong versions: the new line with the old step ($449$ and $229$ failing points), and the new step at $\eta=0.3>\tfrac14$ (FrB1 and FrB2 break).

At $\eta=1/\ell$ the published audits of~\cite{I} are reproduced exactly by the code of this paper, in which $\eta$ and $1/\ell$ are separate arguments (App.~\ref{app:E}).

\subsection{The conclusion}
We searched for states where the lines of this paper give $(F^2-1)\ell<10$, with $F=\min_a\max_k$ of the ratio of the best option to the target (the quantity of Sec.~VIII of~\cite{I}), at the ends A and B in both forms and, for the form with the remainders at $\bar\eta$, at $\lambda=0.9$.
The families are the twelve of~\cite{I} (corner, global, five edge families, five narrow families) and the minimisers of the searches of~\cite{I} as seeds; $n=800$ walkers per family, shrinking random walks and pattern search.
\begin{center}
\small
\begin{tabular}{lccccc}
$\ell$ & A & mid & B & A (chord) & B (chord)\\ \hline
$55.26$ & $22.89$ & $23.66$ & $24.54$ & $22.89$ & $24.54$\\
$960.8$ & $28.22$ & $28.66$ & $29.16$ & $28.22$ & $29.16$\\
$10^4$ & $35.83$ & $36.25$ & $36.73$ & $35.83$ & $36.73$\\
$10^6$ & $51.83$ & $52.26$ & $52.73$ & $51.83$ & $52.73$\\
\end{tabular}
\end{center}
The smallest value, $22.89$, is more than twice what the claim needs.
The forms agree at the minimisers because there $k_E\approx\sqrt r$ is small and the remainders hardly act; the minimum moves from Paper~I's $17.56$ to $22.89$ because the end A has gains at factor one instead of $1/(1+1/\ell)$.
Before the interval proofs, the corner family at $\ell=55.26$ was searched once more with $n=1500$ at both ends: $22.89$ and $24.54$.

Power: dropping a main gain line (P2z, or hFz) makes the margin negative at both ends and both $\ell$ (8 versions); the three planted dips of~\cite{I} (corner, narrow, low $\mu$; depth 5) were found, $11$ of $12$ at once and the last one (corner, end A, $\ell=55.26$) with $n=1500$ or with the widths tripled; the corner dip was found again by the $n=1500$ search at both ends.
A first design of the planted dip (width $0.5$ next to the minimiser) was too narrow to be felt by the walkers at $\ell=55.26$; we replaced it by the design of~\cite{I} and report it here.

\emph{In the lines of this proof, the coefficient cannot be one.}
With the target not scaled ($\gamma$ in place of $(\lambda/\Upsilon)\gamma$) at $\eta=\tfrac14$, the smallest $(F^2-1)\ell$ of the lines is $-26.5$ at $\ell=55.26$ and $-604$ at $\ell=960.8$: $F$ falls to about $3/5$, the sandwich ratio.
This is a statement about the lower bounds of this proof, not about $g$.
The sandwich puts the worst gain factor $\lambda$ and the worst cost factor $\Upsilon$ at the same time; in one window the two need not occur together.
The next subsection asks what the true quantity does.

\subsection{The true quantity at a strong ripple}
\label{sec:truth}
We computed the true quantity $Q=(a)+(a')$, which bounds $g$ through Eq.~\eqref{eq:ga}, on windows with $\eta=\tfrac14$, $0.1$ and $1/\ell$, and searched for small $Q\sqrt\ell$ \numst.

\emph{The method.}
The quantity of \code{dp\_truth.py} of~\cite{I} is a two-dimensional quadrature in $(v,w)$ and works only for moderate $\sigma$.
We write it in the variables of the local lemmas of~\cite{I}.
The integrand is split into a sum of terms (a function of the ripple phase $t$) $\times\,\hat s^k\times\{1,\,1-\cos(\sqrt y\hat s),\,\sin(\sqrt y\hat s)\}$, $k\le2$; the nonlinearity in the ripple, $1/h$ included, stays inside the functions of $t$.
These are sampled on $M$ points and replaced by their discrete Fourier coefficients up to a cut $n_{\rm cut}$; the Gaussian averages $E[\hat s^kU_j(\hat s)e^{in\sqrt\tau\hat s}]$ are then taken in closed form.
For a narrow focus the $\hat s$-average is done instead by $128$-point Gauss--Hermite quadrature; the $w$-average by Gauss--Legendre panels at the two scales $s$ and $L$ (for $L>3$, the part without ripple in the closed form of Lemma~11 of~\cite{I}).
$1-\cos$ is written without cancellation, and everything is in extended precision, so that $\sigma=e^{-960.8}$ is representable.
No term is expanded or dropped; the only errors are those of the quadratures and of the Fourier truncation and aliasing.
Checks: against the quadrature of~\cite{I} at $\ell=8$ ($44$ comparisons, relative difference at most $1.1\times10^{-5}$, the error of that quadrature), against the closed form of Lemma~11 of~\cite{I} at $\ell=8$, $55.26$, $960.8$ ($27$ states, $2\times10^{-15}$), and the grids doubled ($5\times10^{-12}$); every value reported is computed on the finer of the two $t$-grids of the code.

\emph{The windows and states.}
Windows are built from a non-negative unblurred profile with (P2) and (P3), aligned so that $\sup|R|=\eta$ exactly: one harmonic; two harmonics ($m=3,21,61$); a spike (all harmonics up to the blur); a packet of fine harmonics at the trough of the fundamental, where $h=1-\eta$ (the shape that brings the worst cost and the worst gain of the sandwich closest together), and the same packet at the crest.
States: the narrow room (Theorem~\ref{thm:Jnarrow}), the focus at the resolution (Theorem~\ref{thm:Jres}) and the band $8\ell+100\le B<\ell^4$ of Theorem~\ref{thm:Jwide}, all states, and separately the ripple-dominated states (the ripple carries at least half of the denominator).
Search: $400$ points of the unit cube of the parameters, then Nelder--Mead from the best six, for each of the $126$ combinations at $\ell=55.26$; the minimisers re-evaluated on finer grids; the best combinations again at $\ell=200$ and $960.8$.

\emph{The result.}
\begin{center}
\small
\begin{tabular}{lccc}
smallest $Q\sqrt\ell$ found & $\eta=\tfrac14$ & $\eta=0.1$ & $\eta=1/\ell$\\ \hline
$\ell=55.26$ & $3.40$ & $3.49$ & $5.20$\\
$\ell=200$ & $10.3$ & $20.0^*$ & $8.84$\\
$\ell=960.8$ & $22.1$ & $17.1$ & $29.3$\\
\end{tabular}
\end{center}
At $\ell=200$ and $960.8$ only the best combinations of $\ell=55.26$ were searched, with a smaller budget, so these are upper estimates of the minimum (${}^*$: the search ended at a dark ripple).
The smallest values sit in the narrow room, at packets of fine harmonics whose local fineness is $\mu\approx1.2$; there $Q\approx1/(2\sqrt\mu)$, far above $\Gamma_\ell$.
At the resolution the minimum is that of a dark ripple ($7.43$ at $\ell=55.26$, the same for every $\eta$), and in the band of Theorem~\ref{thm:Jwide} every search ends at $Q\sqrt\ell=\sqrt\ell/(8\sigma\ell^4)$, about $10^{17}$ at $\ell=55.26$ and $e^{934.7}$ at $\ell=960.8$ \derivst: there $y=1/(2B)$ is not small, a ripple that carries the denominator pays $(a)+(a')\gtrsim y/(4\sigma)=1/(8\sigma B)$ through the spread of its phase over the focus (the chirp line cKz), and a dark ripple pays the background $C_0/\sqrt b$.
The minima grow roughly like $\sqrt\ell$, since in this class the local fineness $\mu$ of a bright ripple stays near one.
At $\ell=55.26$ a strong ripple gives \emph{smaller} minima than a weak one, because it can be brighter; at $\ell=200$ and $960.8$ the values are upper estimates from a smaller budget, and their order can change with it. The comparison that does not depend on the budget is at the same window shape and the same state: there $Q$ with $\eta=1/\ell$ is $1.4$ to $3.4$ times larger than with $\eta=\tfrac14$.

\emph{Does the sandwich act in the true quantity?}
At each minimiser we also computed $Q_{h1}$, the same window and state with the factors $1/h$ of the scores replaced by one: these factors are what the ripple sandwich bounds by $\lambda$ and $\Upsilon$.
At the ripple-dominated minimisers with $\eta=\tfrac14$, $Q/Q_{h1}$ lies between $0.93$ and $1.04$ (against $(1-\eta)/(1+\eta)=0.6$).
A search that minimises $Q/Q_{h1}$ itself finds $0.71$ at $\eta=\tfrac14$ and $0.87$ at $\eta=0.1$ (against $0.6$ and $0.82$), both for the packet at the trough, as designed; but there $Q\sqrt\ell$ is $257$ and $766$.
So the factors of the sandwich do act in the true quantity, up to most of their bound, but only at states where the information is plentiful; at the states where it is scarce they move $Q$ by a few per cent.

\emph{Power.}
At the overall minimiser and at the ripple-dominated one: leaving out the term (III) of the numerator lowers $Q$ by $43$--$50\,\%$ and the search finds values near zero; the scores computed with $1-ER$ in place of $h$ change $Q$ by up to $15\,\%$ over the minimisers; a planted bowl at a random point, with its bottom at $0.3$ times the minimum, is found by the same search (bottom $1.05$, found $1.06$; and $1.02$, found $1.03$); and an artificial worst sandwich (the brightness in the denominator times $1/(1-\eta)^2$, in the numerator times $1/(1+\eta)^2$) gives $2.11$ and $2.20$ against the true $6.68$ and $3.40$ of the same family: within these families, the search would have found a lowering by the sandwich ratio; it found none.

\emph{What this shows and what it does not.}
No window found has $((a)+(a'))\sqrt\ell$ near one, let alone near $(1-\eta)/(1+\eta)$ (and $g\sqrt\ell\ge(1-\bar\varepsilon)((a)+(a'))\sqrt\ell$).
But in the regime of the class ($x\le0.003$, the fundamental the largest harmonic) the true quantity is far above $\Gamma_\ell$ for weak ripples too.
The states where the law is nearly tight are those of the comb of~\cite{F}, whose period is at the scale of the blur ($x\approx\ell/4$ for its own fundamental); there $(a)+(a')$ is $1.24/\sqrt\ell$ at $\ell=55.26$ and tends to $1/\sqrt\ell$ (\code{dh\_true.py} in the repository of~\cite{I}), and they are outside the class of Theorem~\ref{thm:J}.
So the searches do not test the factor near tightness; they show that, where they look, a strong ripple does not take a fraction of the information away.
\begin{conjecture}[coefficient one for every ripple] \conjst
\label{conj:one}
For every admissible window ($\eta\le\tfrac14$), every state and every $\ell\ge55.26$, $g\sqrt\ell\ge1-\bar\varepsilon$.
\end{conjecture}

\subsection{The convexity inventory}
The parser of Lemma~\ref{lem:chord} was itself tested with three wrong versions of the lines (a loss multiplied by $1-\eta$, $1+\eta$ replaced by $1/(1+\eta)$, $\sqrt\eta$ in d3): all three are rejected.
Our first hand-written inventory had three wrong rows (N4G, B2s, uB2s listed basic functions they do not contain); the parser found them, and Table~\ref{tab:inventory} is the corrected one.
Numerically, every product of up to three of the six functions, and every loss line of~\cite{I} at random states, lies under its chord on a grid of $\lambda$ (smallest difference $\ge-10^{-15}$ relative, rounding).
